\section{Discussion}
\label{sec:discussion}

The discussion separates three claims: whether the contract detects
construction-specific scene defects, whether constrained repair resolves those
defects reliably, and whether improved readiness transfers to runtime task
outcomes. This separation prevents successful simulation videos from being
treated as evidence that every scene requirement is correct. The
defect-consequence study of \Cref{sec:task-outcomes} speaks to the third
claim and, by inversion, to which scene-model features the first two claims
must protect.

\subsection{Which scene-model features are task-critical}
\label{sec:criticality}

The paired defect-injection results yield a criticality ordering that is not
aligned with the effort distribution of current simulation-readiness
specifications, which emphasize asset-level geometry and physical fidelity.

\paragraph{Referential integrity is discrete-fatal.} A single wrong task
reference, a stale construction state, or a wrong semantic label drove task
success from 39\% to exactly zero under both seed partitions. These features
define \emph{what} the executor is asked to do; no amount of perceptual or
control robustness can compensate for executing the wrong target. Scene
contracts should therefore treat reference closure---every task reference
resolves to an existing, correctly labelled, in-place member---as a hard
acceptance gate rather than a soft quality score. Notably, a \emph{missing}
semantic label is far less harmful than a \emph{wrong} one (6.5\% vs.\ 0.0\%):
the executor's nominal-pose prior occasionally suffices, so completeness
requirements can be graded while correctness requirements cannot.

\paragraph{Assembly-interface tolerance is the most sensitive continuous
feature.} Reducing the effective radial clearance of a nominal \O22\,mm hole
from 1\,mm to 0.75\,mm---an error well within what manual modelling or
careless unit handling can introduce---collapsed success from 39\% to 5\%.
The per-episode diagnostics show the defect acting through the intended
physical channel (insertion travel truncated from $\approx$50\,mm to below
5\,mm by the clearance gate), so this is a genuine jamming phenomenon rather
than a metric artefact. Interface dimensions therefore require certified
values with provenance; an interface field without provenance should be
treated as absent.

\paragraph{Geometric accuracy matters relative to the perceptual anchor, not
in absolute terms.} A 10\,mm shift of the entire connection node was almost
fully compensated by closed-loop perception ($-2.5$\,pp, n.s.), whereas the
same 10\,mm between the certified hole position and the observed member cost
5.5\,pp ($p<0.001$). The scene model's coordinate guarantees should
accordingly be specified and verified \emph{relative to the features the
controller can observe}, since absolute-frame accuracy is partially
self-correcting at runtime while anchor-relative error is invisible to the
controller by construction.

\paragraph{Physical-property criticality is stage-dependent.} Bolt mass
perturbations of $\times0.5$ and $\times2$ produced bit-identical outcomes
because the transport-and-insert executor holds the bolt in a kinematic
virtual grasp. We read this as a structural property of the motion stage, not
evidence that physical fidelity is unimportant: the same defect is expected to
become consequential in any stage with forceful contact, and a full-contact
ablation with a physically grasped bolt is required before the physical family
can be ranked.

\subsection{Sensitivity and robustness of the findings}
\label{sec:sensitivity}

Three analyses bound the statistical fragility of these conclusions.

First, all fourteen conditions were re-run under an independent seed
partition (seed 123, 200 paired episodes per condition). Success rates
deviate by at most 2.5\,pp from the seed-42 values; every condition that was
significant remains significant with the same direction and comparable
magnitude, and every non-significant condition remains non-significant. The
criticality ordering is therefore a property of the scene-model features, not
of a particular random draw of episodes.

Second, the severity gradients are monotone where the defect acts through a
continuous channel. Hole-position error at $n=1000$ paired episodes gives
37.5\%, 36.8\%, and 31.3\% success for 2, 5, and 10\,mm against a 38.1\%
baseline; interface clearance gives 5.0\%, 4.0\%, and 2.5--2.0\% for 0.75,
0.5, and 0.25\,mm across the two seeds. Monotone dose--response supports a
causal interpretation of the injected defect rather than threshold noise in
the success criterion.

Third, the power upgrade from 200 to 1000 paired episodes changed exactly one
conclusion and clarifies its meaning. At $n=200$ the 2\,mm and 5\,mm
hole-position defects were underpowered ($b$=1 and 4 discordant pairs); at
$n=1000$, 5\,mm is significant ($p=0.015$, corroborated at $n=200$ under
seed 123 with $p=0.002$) while 2\,mm remains non-significant ($p=0.109$).
For this executor, hole-position errors below roughly 5\,mm lie inside the
policy's tolerance band, and claims about sub-5\,mm criticality would require
either a stricter executor or substantially larger samples. We accordingly
report effect sizes alongside significance and do not rank features on
$p$-values alone.

\subsection{Boundary of pure-simulation conclusions}
\label{sec:boundary}

All runtime evidence in this study is simulated, and the conclusions should
be read as structural rather than numerical. Rigid-body simplification,
approximate contact parameters, and simplified pins in place of threaded
fasteners mean that absolute success rates and clearance thresholds will
shift on hardware; what transfers is the ordering and the causal channel
(identification of \emph{which} defect classes are fatal, graded, or
self-correcting). Three further boundaries apply. (i) A single
robot--controller combination was used: the compensation of whole-node pose
error depends on the executor's $\pm$5\,cm pose domain randomization during
training, and an executor without such randomization would likely show a
larger member-pose effect. (ii) The physical-family null result is an
artefact of the virtual-grasp scaffold and bounds only the
transport-and-insert stage. (iii) The success criterion (12\,mm alignment,
20\,mm insertion depth) matches the training reward; stricter assembly
tolerances would shift all conditions downward but preserve pairing. A
physical-robot validation study is deferred to future work, and the present
results should be treated as lower bounds on defect consequence: defects that
are benign here are not certified benign on hardware, while defects that are
fatal here are fatal by construction of the task definition.

\subsection{Implications for cross-runtime specifications}
\label{sec:cross-runtime}

The defect-consequence profile suggests which validation rules belong in a
runtime-agnostic scene contract and which require per-runtime calibration.
Reference closure, semantic-identity correctness, and construction-state
consistency are fatal regardless of executor and should be specified once, as
cross-runtime rules with hard-fail semantics. Interface-tolerance and
anchor-relative-geometric rules have executor-dependent thresholds (here
$\approx$1\,mm clearance and $\approx$5\,mm anchor-relative position) and
should be specified parametrically, with the threshold supplied by the target
controller's tolerance band and verified by the kind of paired injection
benchmark used here. Asset-level physical fidelity, finally, should be
required only for task stages with a dynamical channel; mandating it
uniformly imposes authoring cost without measurable task benefit.
